% ============================================================================
% Section 8.1 — Populations and Samples
% CCSS 7.SP.A.1
% ============================================================================
\section{Populations and Samples}

\begin{stepGoal}
\begin{itemize}
    \item Tell the difference between a population and a sample
    \item Decide whether a sample is random and representative
\end{itemize}
\end{stepGoal}

\begin{stepVocabBox}
    \vocabItem{Population}{The \textbf{entire} group you want to learn about.}
    \vocabItem{Sample}{A \textbf{smaller} group chosen from the population.}
    \vocabItem{Random Sample}{Every member of the population has an \textbf{equal chance} of being chosen.}
    \vocabItem{Bias}{When a sample does \textbf{not} represent the population fairly.}
\end{stepVocabBox}

\begin{stepsCard}[How to Identify and Evaluate a Sample]
    \stepItem{Identify the \textbf{population} — the whole group you want to study.}
    \stepItem{Identify the \textbf{sample} — the smaller group actually surveyed or measured.}
    \stepItem{Check: Was the sample chosen \textbf{randomly}? Does every member have an equal chance?}
    \stepItem{Decide if the sample is \textbf{representative} — if not, explain the bias.}
\end{stepsCard}

% --- Worked Example 1 ---
\begin{stepExample}{A city has $20{,}000$ residents. A reporter surveys $150$ people at a shopping mall about a new park. Is this a good sample?}
    \stepShow{1} Population: all $20{,}000$ residents of the city.

    \stepShow{2} Sample: the $150$ people surveyed at the mall.

    \stepShow{3} Not random — only people who happen to be at the mall are included.

    \stepShow{4} Biased — mall shoppers may not represent all residents (e.g., children and elderly may be missing).
    \stepResult{The sample is biased — a random method like calling randomly chosen phone numbers would be better.}
\end{stepExample}

% --- Worked Example 2 ---
\begin{stepExample}{A teacher puts all $30$ student names in a hat and draws $10$. She asks them about homework time. Evaluate the sample.}
    \stepShow{1} Population: all $30$ students in the class.

    \stepShow{2} Sample: the $10$ names drawn.

    \stepShow{3} Random — every name had an equal chance of being drawn.

    \stepShow{4} Representative — since the selection is random, it should reflect the class fairly.
    \stepResult{This is a good random sample.}
\end{stepExample}

\watchOut{Surveying only your friends or only people in one location is almost always biased, even if the sample is large!}

% ============================================================================
% PRACTICE
% ============================================================================
\newpage
\begin{stepPractice}{Populations and Samples Practice}
    \resetProblems

    \stepPracticeHeader[step1Color]{Identify the Population}
    \prob A principal surveys $50$ of $400$ students about school lunches. What is the population? \answerBlank[3cm]
    \answer{All $400$ students}

    \stepPracticeHeader[step2Color]{Identify the Sample}
    \prob In the same survey, what is the sample? \answerBlank[3cm]
    \answer{The $50$ students surveyed}

    \stepPracticeHeader[step3Color]{Is the Sample Random?}
    \prob A store asks every $10$th customer to rate service. Is this random? \answerBlank[3cm]
    \answer{Yes — systematic sampling gives each customer an equal chance}

    \prob A coach surveys only the starting five players about practice length. Is this random? \answerBlank[3cm]
    \answer{No — only starters are chosen, bench players are left out}

    \stepPracticeHeader[step4Color]{Evaluate the Sample}
    \wordProblem{A TV station polls viewers who call in about a new law. Is this a representative sample of the whole city? Explain.}{~}
    \answer{No — only people who feel strongly enough to call are included, so the sample is biased.}
\end{stepPractice}
